\section{黎曼度量}
\label{sec:riem-metric}

切空间$T_pM$是线性空间，可以把切向量相加、数乘，却还不能回答“这个速度有多快”“这两个方向是否垂直”。欧氏空间中这两个问题由内积回答，流形上自然的做法是给每个切空间各配一个内积。困难在于不同点的切空间是不同的线性空间，内积必须逐点指定；为了能对长度求积分、对内积求导，还要求它随位置光滑变化。

\paragraph{度量的定义}
设$M$是$n$维光滑流形。在坐标图$(U,\varphi)$中记坐标为$x=(x^1,\ldots,x^n)$，上标是坐标编号而不是幂；相应的坐标切基记为$\partial_1,\ldots,\partial_n$，它们在每点$p\in U$构成$T_pM$的基（第~\ref{chap:manifold-basics}章）。切向量$u\in T_pM$写成$u=\sum_iu^i\partial_i$，数组$(u^1,\ldots,u^n)$是它在这组基下的分量。

\begin{definition}[黎曼度量]
\label{def:riem-metric}
光滑流形$M$上的\term{黎曼度量}（Riemannian metric）$g$为每个点$p\in M$指定切空间上的一个内积
\[
g_p:T_pM\times T_pM\to\mathbb R,
\]
即对称、双线性且正定的函数，并且满足光滑性：在任一坐标图中，函数$g_{i,j}(x)=g_p(\partial_i,\partial_j)$（$p=\varphi^{-1}(x)$）都是$x$的光滑函数。带有黎曼度量的光滑流形$(M,g)$称为\term{黎曼流形}（Riemannian manifold）。
\end{definition}

定义里的三项要求各有分工。对称与双线性使$g_p$由矩阵$\bm G(x)=[g_{i,j}(x)]$完全确定：若$u,v$的分量为$(u^i)$与$(v^j)$，则
\begin{equation}
g_p(u,v)=\sum_{i,j=1}^n g_{i,j}(x)\,u^iv^j .
\label{eq:riem-metric-coordinates}
\end{equation}
正定性要求$\bm G(x)$对称正定，于是长度$\lVert u\rVert_g=\sqrt{g_p(u,u)}$只在$u=0$时为零，非零向量的夹角$\vartheta$由$\cos\vartheta=g_p(u,v)/(\lVert u\rVert_g\lVert v\rVert_g)$确定。光滑性则把逐点的内积连成一个整体对象；后面求联络系数时要对$g_{i,j}$求导，这一条件不可省。

光滑性条件中的“任一坐标图”并不需要逐张核对。换用另一组坐标$y$时，记$\bm A(y)=[\partial x^i/\partial y^a]$为坐标变换的Jacobian矩阵。同一个切向量的旧分量等于$\bm A$乘新分量，代入式\eqref{eq:riem-metric-coordinates}得到
\begin{equation}
\widetilde{\bm G}(y)=\bm A(y)^{\top}\bm G(x(y))\bm A(y).
\label{eq:riem-metric-transform}
\end{equation}
右端各因子都光滑，所以一张图中光滑，重叠处的另一张图中也光滑。这个变换规律也说明，$\bm G$的条目本身随坐标改变，不变的是它所规定的内积。

\paragraph{同一几何的不同坐标}
欧氏平面取标准内积。直角坐标中$\bm G=\bm I$；改用极坐标$(r,\theta)$，$r>0$，坐标切基为$\partial_r=(\cos\theta,\sin\theta)^{\top}$与$\partial_\theta=(-r\sin\theta,r\cos\theta)^{\top}$，于是
\[
\bm G(r,\theta)=\begin{bmatrix}1&0\\0&r^2\end{bmatrix}.
\]
系数随位置变化，并不说明几何改变了：$\partial_\theta$表示“角度增加一单位”的运动，在半径$r$处走过的弧长约为$r$。度量矩阵记录的是坐标单位与真实长度之间的换算，同一几何在不同坐标中可以有不同的矩阵。

\paragraph{嵌入诱导的度量}
曲面通常坐落在欧氏空间中，此时最自然的长度就是环境中的长度。设$M\subseteq\mathbb R^D$是嵌入子流形（第~\ref{chap:manifold-relations}章），局部参数化$\Phi:V\to M$的Jacobian矩阵$\bm J(x)=[\partial\Phi/\partial x^1,\ldots,\partial\Phi/\partial x^n]$满列秩，它的各列正是坐标切向量在$\mathbb R^D$中的表示。把环境内积限制到切空间，得到
\begin{equation}
\bm G(x)=\bm J(x)^{\top}\bm J(x),\qquad g_{i,j}=\Bigl\langle\frac{\partial\Phi}{\partial x^i},\frac{\partial\Phi}{\partial x^j}\Bigr\rangle .
\label{eq:riem-induced-metric}
\end{equation}
满列秩保证$\bm a^{\top}\bm G\bm a=\lVert\bm J\bm a\rVert_2^2>0$对非零$\bm a\in\mathbb R^n$成立，所以$\bm G$正定。这样得到的度量称为\term{诱导度量}（induced metric）。若参数化在某点秩亏，同一公式只给出半正定矩阵，不再满足黎曼度量的定义。

\begin{example}[单位球面的诱导度量]
\label{ex:riem-sphere-metric}
单位球面$\mathbb S^2\subseteq\mathbb R^3$在去掉两极和一条经线后，可用极角$\theta\in(0,\pi)$与方位角$\phi\in(0,2\pi)$参数化：
\[
\Phi(\theta,\phi)=(\sin\theta\cos\phi,\ \sin\theta\sin\phi,\ \cos\theta)^{\top}.
\]
两列切向量为$\partial_\theta\Phi=(\cos\theta\cos\phi,\cos\theta\sin\phi,-\sin\theta)^{\top}$与$\partial_\phi\Phi=(-\sin\theta\sin\phi,\sin\theta\cos\phi,0)^{\top}$，二者内积为零，长度平方分别为$1$与$\sin^2\theta$，所以
\[
\bm G(\theta,\phi)=\begin{bmatrix}1&0\\0&\sin^2\theta\end{bmatrix}.
\]
沿纬线（固定$\theta$）方位角增加$\Delta\phi$，走过的长度是$\sin\theta\,\Delta\phi$：越靠近两极，同样的角度变化对应越短的路程。
\end{example}

\paragraph{同一流形上的不同度量}
诱导度量只是一种选择。同一个光滑流形可以配上许多不同的黎曼度量，它们给出不同的长度和角度。例如在上半平面$H=\{(x,y):y>0\}$上，除了欧氏度量$\bm G=\bm I$，还可以取
\[
\bm G(x,y)=\frac1{y^2}\begin{bmatrix}1&0\\0&1\end{bmatrix}.
\]
两者在每点只差一个正的倍数，所以角度相同；长度却不同：在这种度量下，竖直线段从$y=1$走到$y=\varepsilon$的长度是$\int_\varepsilon^1 y^{-1}\,\mathrm dy=\ln(1/\varepsilon)$，当$\varepsilon\to0$时趋于无穷，而欧氏长度不超过$1$。下一节的距离、第~\ref{chap:manifold-calculus}章的体积和梯度都由度量决定，所以凡是谈论这些量，都要说明所用的是哪一个度量。

\paragraph{度量的存在}
一般光滑流形不一定嵌在给定的欧氏空间中，但总存在黎曼度量。取一族坐标图$\{(U_\alpha,\varphi_\alpha)\}$和从属于它的光滑单位分解$\{\psi_\alpha\}$（第~\ref{chap:manifold-basics}章），在每张图中先用坐标的欧氏内积定义$g^{(\alpha)}$，再令
\[
g_p(u,v)=\sum_\alpha\psi_\alpha(p)\,g^{(\alpha)}_p(u,v).
\]
$\psi_\alpha$的支撑含于$U_\alpha$，所以每一项在$U_\alpha$外按零延拓后仍光滑；每点附近只有有限项非零，所以和是光滑的；对称与双线性逐项保持。对非零$u$，每项$\psi_\alpha(p)g^{(\alpha)}_p(u,u)\geq0$，且至少有一个$\psi_\alpha(p)>0$，相应一项严格为正，故$g_p$正定。这个构造依赖的是正定矩阵的非负组合仍正定；它只说明度量存在，并不挑出某个“正确”的度量。
